% =====================================================================
%  CHAPTER 1E: 矿物质 — 参考：0310营养参考_extracted.txt (PPT原文)
% =====================================================================
\chapter{矿物质}
\chapterintro{掌握矿物质的定义与分类（常量/微量元素）；掌握钙的分布、功能、吸收影响因素和食物来源（难点*）；掌握铁的功能、缺乏三阶段及血红素铁/非血红素铁吸收特点；掌握碘的分布与来源；熟悉锌/硒的功能。}

\section{矿物质概述}

\begin{defbox}
\textbf{矿物质（Minerals）：}除碳、氢、氧和氮主要以有机化合物形式存在外，人体组织中其余元素统称为矿物质（无机盐/灰分）。

\textbf{分类：}\imp{常量元素（Macroelement）}：占体重0.01\%以上——Ca/P/K/Na/Cl/Mg/S。\imp{微量元素（Microelement）}：占体重0.01\%以下。必需微量元素：Fe/Cu/Zn/Se/Cr/I/Co/Mo。可能必需：Mn/Si/Ni/B/V。潜在毒性：F/Pb/Cd/Hg/As/Al/Sn/Li。

\textbf{营养特点：}(1)除食物外是唯一可通过天然水获取的营养素；(2)体内分布极不均匀；(3)\imp{生理剂量与中毒剂量范围较窄}；(4)功能：构成组织/维持渗透压与酸碱平衡/维持神经肌肉兴奋性及细胞膜通透性/构成活性物质/酶激活剂。
\end{defbox}

\section{钙（Calcium）}

\begin{keybox}
\textbf{分布：}\imp{人体含量最多的无机元素}（1.5\%$\sim$2.0\%），总量约1200g，\imp{99\%在骨骼牙齿中}（羟磷灰石Ca$_{10}$[PO$_4$]$_6$[OH]$_2$），1\%为混溶钙池。

\textbf{功能：}构成骨骼牙齿；酶激活剂/参与细胞代谢；参与凝血；构成混溶钙池；维持神经肌肉应激性；维持心跳节律。

\textbf{吸收影响因素：}\imp{促进}——乳糖/氨基酸/VD；\imp{抑制}——植酸盐/膳食纤维/草酸盐/脂肪/乙醇。\textbf{来源：}小虾皮/奶与奶制品/豆与豆制品。
\end{keybox}

\section{铁（Iron）}

\begin{keybox}
\textbf{功能：}血红蛋白原料（O$_2$/CO$_2$运输）；多种酶原料（生物氧化）；促进胡萝卜素→VA；促进抗体产生/胶原合成。

\textbf{缺乏：}好发人群——婴幼儿/青少年/育龄妇女。\imp{铁缺乏三阶段：}铁减少期→缺铁性红细胞生成期→缺铁性贫血期。

\textbf{吸收：}\imp{血红素铁}：与卟啉结合，直接吸收，不受干扰。\imp{非血红素铁}：以Fe(OH)$_3$形式存在于植物性食物中，须经胃酸分离并还原为二价铁才可吸收，吸收率低且受干扰。
\end{keybox}

\section{磷与镁}

\begin{defbox}
\textbf{磷（Phosphorus）：}
\begin{itemize}[leftmargin=*]
    \item 人体含量约650g，\imp{约70\%在骨骼中}，30\%在软组织
    \item 功能：构成骨骼牙齿（羟磷灰石）；核酸/磷脂/含磷辅酶（NAD$^+$、NADP$^+$、TPP等）的组成成分；参与能量代谢（ATP）；维持酸碱平衡
    \item 吸收：受维生素D促进；植酸抑制
    \item 钙磷比（Ca:P）：成人适宜1:1$\sim$1:1.5，婴儿母乳2:1（最适宜）
\end{itemize}

\textbf{镁（Magnesium）：}
\begin{itemize}[leftmargin=*]
    \item 成人含镁约25g，\imp{60\%$\sim$65\%在骨骼中}，27\%在肌肉
    \item 功能：300多种酶的辅因子；参与蛋白质合成和能量代谢；维持神经肌肉兴奋性；心血管保护作用
    \item 缺乏：神经肌肉兴奋性增高、心律失常、低钙血症（镁缺乏同时导致低血钙）
    \item 来源：绿叶蔬菜（叶绿素含镁）、坚果、豆类、全谷物
\end{itemize}
\end{defbox}

\section{锌（Zinc）}

\begin{keybox}
\textbf{功能：}
\begin{itemize}[leftmargin=*]
    \item 多种酶的组成成分和激活剂（超200种含锌酶）：碱性磷酸酶、RNA/DNA聚合酶、碳酸酐酶
    \item 促进生长发育和组织再生：参与核酸和蛋白质合成
    \item 维持正常味觉和食欲：味觉素的组成成分
    \item 参与免疫功能：促进淋巴细胞增殖和抗体生成
    \item 抗氧化：超氧化物歧化酶（SOD）的组成成分
    \item 促进维生素A代谢：参与视黄醇结合蛋白合成
\end{itemize}

\textbf{缺乏：}婴幼儿/儿童——\imp{生长发育迟缓、性发育不全、脑发育受损、食欲不振、异食癖}；成人——味觉减退、伤口愈合延迟、免疫功能下降、脱发、皮肤损害。

\textbf{来源：}动物性食物（贝壳类海产品、红肉、动物内脏）锌含量和利用率均高；\imp{母乳锌含量>牛乳}，初乳最高，生物价值更好；植物性食物受植酸抑制吸收。

\textbf{促进吸收因素：}动物蛋白质、组氨酸、半胱氨酸、柠檬酸。\textbf{抑制因素：}植酸、膳食纤维、高钙、高铜、高镉。
\end{keybox}

\section{硒（Selenium）}

\begin{keybox}
\textbf{功能：}
\begin{itemize}[leftmargin=*]
    \item \imp{谷胱甘肽过氧化物酶（GSH-Px）}的重要组成成分——抗氧化、保护细胞膜
    \item 参与甲状腺激素代谢：碘甲腺原氨酸脱碘酶（将T$_4$转化为活性T$_3$）
    \item 维持免疫功能：促进抗体生成
    \item 重金属解毒：与汞、镉、铅等结合形成不溶性硒化物
    \item 保护心血管、抗肿瘤
\end{itemize}

\textbf{缺乏：}\imp{克山病（Keshan Disease）}——以心肌坏死为特征的地方性心肌病（中国东北至西南的克山病带，硒缺乏为主因）；\imp{大骨节病（Kaschin-Beck Disease）}——骨关节变形、软骨坏死。过量可引起\imp{硒中毒}（脱发、脱甲、神经系统损害）。

\textbf{来源：}动物肝、肾、海产品、肉类；植物性食物硒含量取决于土壤硒水平。中国成人RNI：60 $\mu$g/d。
\end{keybox}

\section{碘（Iodine）}

\begin{defbox}
人体含碘约20$\sim$50mg，\imp{70\%$\sim$80\%在甲状腺}。

\textbf{功能：}参与甲状腺激素（T$_3$、T$_4$）合成，调节基础代谢、促进生长发育（尤其是神经系统发育）。

\textbf{缺乏：}\imp{地方性甲状腺肿（Goiter）}——缺碘→T$_3$/T$_4$合成减少→TSH反馈性升高→甲状腺代偿性肿大。\imp{克汀病（Cretinism）}——孕期/婴幼儿期严重缺碘导致不可逆的智力低下和生长发育障碍。缺碘地区：食用碘盐（食盐加碘是最经济有效的预防措施）。

\textbf{来源：}海盐/海产品（海带/紫菜/海鱼/贝类等最丰富）。成人RNI：120 $\mu$g/d。
\end{defbox}

\section{其他微量元素}

\begin{defbox}
\textbf{铜（Copper）：}参与铁代谢（铜蓝蛋白促进铁的吸收和利用）；参与结缔组织中胶原蛋白和弹性蛋白的交联（赖氨酰氧化酶）；参与黑色素合成（酪氨酸酶）。缺乏：小细胞低色素性贫血（与铁代谢障碍有关）、中性粒细胞减少、骨质疏松。

\textbf{铬（Chromium）：}\imp{三价铬是葡萄糖耐量因子（GTF）的组成成分}，增强胰岛素作用，促进葡萄糖利用。缺乏：葡萄糖耐量降低、胰岛素抵抗。来源：全谷物、肉类、啤酒酵母。成人AI：30 $\mu$g/d。

\textbf{氟（Fluoride）：}促进骨骼和牙齿的矿化，预防龋齿。\imp{缺乏→龋齿发病率增高；过量→氟斑牙（Dental Fluorosis）和氟骨症（Skeletal Fluorosis）}。来源：饮用水（主要来源）、海产品、茶叶（含氟量高）。生理剂量与中毒剂量范围极窄！

\textbf{钴（Cobalt）：}维生素B$_{12}$的组成成分（唯一含金属元素的维生素），通过VitB$_{12}$发挥生理作用。缺乏即VitB$_{12}$缺乏——巨幼红细胞性贫血。
\end{defbox}

\section{复习题}
\begin{enumerate}[leftmargin=*, itemsep=8pt]
    \item \textbf{【名词解释】}矿物质；常量元素；微量元素；混溶钙池；血红素铁；非血红素铁；克山病；克汀病；葡萄糖耐量因子（GTF）
    \item \textbf{【简答题】}简述矿物质的分类和营养特点。
    \item \textbf{【简答题】}简述钙的分布、生理功能、吸收影响因素和食物来源。
    \item \textbf{【简答题】}简述铁的功能、缺乏三阶段及血红素铁与非血红素铁的区别。
    \item \textbf{【简答题】}简述锌的生理功能及缺乏表现。
    \item \textbf{【简答题】}简述硒的生理功能、缺乏病（克山病/大骨节病）及食物来源。
    \item \textbf{【简答题】}简述碘的分布、功能及缺乏病（地方性甲状腺肿/克汀病）。
    \item \textbf{【简答题】}简述钙磷比的意义及磷、镁的主要功能。
\end{enumerate}

\subsection*{参考答案要点}
\begin{enumerate}[leftmargin=*, itemsep=5pt]
    \item 见正文定义。\textbf{克山病}：以心肌坏死为特征的地方性心肌病，硒缺乏为主因。\textbf{克汀病}：孕期/婴幼儿期严重缺碘导致的不可逆智力低下和生长发育障碍。\textbf{GTF}：三价铬组成的葡萄糖耐量因子，增强胰岛素作用。
    \item 分类：常量(>0.01\%：Ca/P/K/Na/Cl/Mg/S)+微量(<0.01\%：必需Fe/Cu/Zn/Se/Cr/I/Co/Mo；可能必需Mn/Si/Ni/B/V；潜在毒性F/Pb/Cd/Hg/As/Al/Sn/Li)。特点：(1)唯一可通过天然水获取；(2)分布极不均匀；(3)生理剂量与中毒剂量范围窄；(4)构成组织/渗透压酸碱平衡/神经肌肉兴奋性/酶激活剂。
    \item 含量最多无机元素(1.5-2.0\%)，99\%骨骼牙齿(羟磷灰石)，1\%混溶钙池。功能：骨骼牙齿/酶激活/凝血/神经肌肉应激性/心跳节律。促进吸收：乳糖/氨基酸/VD。抑制：植酸盐/草酸盐/膳食纤维/脂肪/乙醇。来源：小虾皮/奶制品/豆制品。成人RNI 800mg/d。
    \item 功能：血红蛋白(O$_2$/CO$_2$运输)/生物氧化/促进胡萝卜素→VA/抗体+胶原。缺乏三阶段：铁减少期→缺铁性红细胞生成期→缺铁性贫血期。血红素铁(与卟啉结合，直接吸收不受干扰)vs非血红素铁(Fe(OH)$_3$形式，需还原为Fe$^{2+}$，吸收率低受植酸/草酸干扰)。男性RNI 12mg/d，女性20mg/d。
    \item 功能：多种酶(超200种)组成成分和激活剂/促进生长发育和组织再生/维持味觉食欲(味觉素)/免疫功能/SOD抗氧化/促进VA代谢。缺乏：生长发育迟缓/性发育不全/脑发育受损/食欲不振/异食癖/味觉减退/伤口愈合延迟/免疫功能下降。来源：贝壳类海产品/红肉/动物内脏，母乳>牛乳。成人男性RNI 12.5mg/d，女性7.5mg/d。
    \item 功能：GSH-Px组成成分(抗氧化/保护细胞膜)/参与甲状腺激素代谢(脱碘酶T$_4$→T$_3$)/维持免疫/重金属解毒。缺乏：克山病(心肌坏死，缺硒地带)、大骨节病(骨关节变形)。过量：硒中毒(脱发/脱甲/神经损害)。来源：动物肝肾/海产品/肉类(植物受土壤硒水平影响)。成人RNI 60$\mu$g/d。
    \item 70-80\%在甲状腺。功能：合成T$_3$/T$_4$，调节基础代谢/促进生长发育(尤其神经系统)。缺乏：地方性甲状腺肿(缺碘→TSH升高→甲状腺代偿性肿大)/克汀病(孕期/婴幼儿缺碘→不可逆智力低下)。来源：海产品(海带/紫菜/海鱼最丰富)。成人RNI 120$\mu$g/d。食盐加碘最经济有效。
    \item Ca:P成人1:1-1:1.5，婴儿母乳2:1(最适宜)。磷功能：骨骼牙齿/核酸磷脂/含磷辅酶/能量代谢(ATP)/酸碱平衡。镁功能：300多种酶辅因子/蛋白质合成/能量代谢/神经肌肉兴奋性/心血管保护。缺乏→神经肌肉兴奋性增高/心律失常/低钙血症。来源：绿叶蔬菜(叶绿素含镁)/坚果/豆类/全谷物。
\end{enumerate}
\newpage
